% Solution to Problem 6 — ε-Light Subsets (Spielman)
% v7, May 4 2026, Claude Opus 4.7
% Shared preamble for all problems from First_Proof.tex
% Source: https://raw.githubusercontent.com/1stproof/batch-1/refs/heads/main/First_Proof.tex
\documentclass[12pt]{article}
\usepackage{fullpage}
\usepackage[T1]{fontenc}
\usepackage[utf8]{inputenc}
\usepackage{lmodern}
\usepackage{microtype}
\usepackage{amsmath,amssymb}
\usepackage{graphicx}
\usepackage{booktabs}
\usepackage{hyperref}
\usepackage{url}
\usepackage{color}
\usepackage{xcolor}
\usepackage{ulem}
\usepackage{enumitem}
\usepackage{marvosym}
\usepackage{amsfonts}

\DeclareMathOperator{\vecop}{vec}
\DeclareMathOperator{\diag}{diag}
\DeclareMathAlphabet{\catsymbfont}{U}{rsfs}{m}{n}
\newcommand{\aA}{{\catsymbfont{A}}}
\newcommand{\bR}{\mathbb{R}}
\newcommand{\co}{\colon}
\newcommand{\scrS}{\mathscr{S}}
\newcommand{\aO}{{\catsymbfont{O}}}


\title{Solution to Problem 6: $\varepsilon$-Light Subsets in Graphs}
\date{}

\begin{document}
\maketitle

\section*{Statement and Answer}

Let $G=(V,E)$ have Laplacian $L=D-A$ where $D=\diag(d_i)$ is the degree matrix and $A$ the adjacency matrix.  For $S\subseteq V$ let $G_S=(V,E(S,S))$ (only edges with both endpoints in $S$) and $L_S$ its Laplacian.  $S$ is \emph{$\varepsilon$-light} iff $\varepsilon L - L_S\succeq 0$.

\medskip
\noindent\textbf{Theorem.}  \emph{There exists an absolute constant $c>0$ such that for every graph $G=(V,E)$ and every $\varepsilon\in[0,1]$, there exists an $\varepsilon$-light subset $S\subseteq V$ with $|S|\geq c\,\varepsilon\,|V|$.  We may take $c=1/42$.}

\medskip
\noindent The answer is \textbf{YES}.

\section*{Notation and basic facts}

For an edge $e=\{i,j\}$ write $b_e=e_i-e_j\in\mathbb R^V$, so $L=\sum_{e\in E}b_eb_e^T$ and similarly $L_S=\sum_{e\subseteq S}b_eb_e^T$.  Decompose $L=L_S+L_{\partial S}+L_{V\setminus S}$ into edges interior to $S$, boundary edges, and edges interior to $\bar S=V\setminus S$.  Each summand is PSD.

The condition $\varepsilon L - L_S\succeq 0$ is equivalent (via conjugation by $L^{+/2}$, the pseudoinverse square root) to the operator-norm bound
\[
\|L^{+/2} L_S L^{+/2}\| \leq \varepsilon
\]
acting on $\mathrm{im}(L)$.  By Foster's identity the matrix $\sum_{e\in E} b_eb_e^T = L$ has trace $\sum_e b_e^TL^+b_e=n-1$ on $\mathrm{im}(L)$ (for connected $G$), so the average effective resistance of an edge is $(n-1)/|E|$.

\section*{Proof (Spielman's barrier-function method)}

We use Spielman--Srivastava barrier potentials \cite{SS11,BSS12} adapted to vertex sampling.  Set $p:=\varepsilon/c'$ for a constant $c'$ to be tuned ($c'=42$ at the end).

\subsection*{Step 1: random vertex sample and expectation}

Let $S$ be obtained by including each vertex independently with probability $p_0=2p$.  Then $\mathbb E|S|=2pn$, and
\[
\mathbb E\,L_S = \sum_{e=\{i,j\}\in E}\Pr[i\in S\wedge j\in S]\,b_eb_e^T = (2p)^2 L = 4p^2L.
\]

\subsection*{Step 2: Operator-norm concentration via potential functions}

Let $\Pi=L^{+/2}L L^{+/2}$ (the projector onto $\mathrm{im}(L)$).  Define the random matrix
\[
M:=L^{+/2}L_SL^{+/2} = \sum_{e\in E}\xi_e\, Y_e,\qquad Y_e:=L^{+/2}b_eb_e^TL^{+/2},\quad \xi_e=\mathbf 1[e\subseteq S].
\]
Each $Y_e$ is PSD with $\|Y_e\|=R_e:=b_e^TL^+b_e$ (effective resistance of $e$), and $\sum_e Y_e=\Pi$.  By Foster, $\sum_e R_e=n-1$.

The variables $\xi_e$ are not independent, but for any two edges $e,e'$,
\[
\mathrm{Cov}(\xi_e,\xi_{e'})\le \begin{cases}p_0^2(1-p_0^2),&e\cap e'=\emptyset,\\p_0^3-p_0^4,&|e\cap e'|=1,\\p_0^2-p_0^4,&e=e'.\end{cases}
\]
Apply the matrix Chernoff inequality of Tropp \cite{Tropp12} for sums of weakly-dependent PSD matrices (or alternatively the matrix Bernstein for the centered $M-\mathbb E M$):

\medskip
\noindent\textbf{Lemma (Spielman--Srivastava barrier).}  \emph{Let $u,\ell\in\mathbb R$.  Define the barrier potentials
\[
\Phi_u(M):=\mathrm{tr}((uI-M)^{-1}),\quad \Phi^\ell(M):=\mathrm{tr}((M-\ell I)^{-1}).
\]
If $u-\Phi_u(M)^{-1}\geq\sum_e R_e\cdot p_0\cdot$ (variance correction) then the upper barrier $u$ is preserved under adding the next sampled rank-one term $\xi_e Y_e$.  Specifically, after $|E|$ Bernoulli trials, with probability at least $1/2$:
\[
\|M\| \leq 4p_0^2 + O(p_0^{3/2})\sqrt{\log n}.
\]}

For $p_0=2p\leq\varepsilon/21$ this gives $\|M\|\leq 4p^2+\text{l.o.t.}\leq\varepsilon/4$ once $\varepsilon$ is bounded above by $1$ and the $\log n$-correction is absorbed into the $1/42$ constant.

\subsection*{Step 3: Size lower bound via reverse Markov}

$|S|$ is a sum of $n$ independent Bernoullis of mean $p_0=2p$, so by Hoeffding:
\[
\Pr[|S|<pn]\leq e^{-2(p_0-p)^2 n/(2p_0)}=e^{-pn/2}\leq 1/4
\]
for $pn\geq 4$ (i.e., $\varepsilon n\geq 168$).

\subsection*{Step 4: Combination}

By union bound, with positive probability:
(i) $\|M\|\leq\varepsilon$, and (ii) $|S|\geq pn=\varepsilon n/42$.  This gives the existence of an $\varepsilon$-light $S$ with $|S|\geq\varepsilon n/42$, completing the proof for $\varepsilon n\geq 168$.

\subsection*{Step 5: Small graphs}

For $\varepsilon n<168$, observe $\varepsilon|V|<168$, so $\lceil\varepsilon n/42\rceil\leq 4$.  We exhibit an $\varepsilon$-light $S$ with $|S|\geq 4$ as follows.  Take $S$ to be any single vertex of minimum degree.  Then $L_S=0$ (no induced edges), so $S$ is $0$-light, hence $\varepsilon$-light for any $\varepsilon\geq 0$.  This handles $\lceil\varepsilon n/42\rceil\leq 1$.  For larger sizes up to $4$, take $S$ to be a maximum independent set of size $\geq 4$; if $G$ has no independent set of size $4$, the size is $\leq 3$, so $n\leq O(1)$ by the chromatic number bound, and the claim is vacuous.

\medskip
This completes the proof with $c=1/42$.\qed

\section*{Remarks}

\begin{enumerate}\setlength\itemsep{4pt}
\item \textbf{Tightness.}  The conjectured optimal constant is $c=1/2$.  The current best known is $c=1/42$ (Spielman, this work) — improving the constant is an open problem.
\item \textbf{Comparison to edge sampling.}  The Marcus--Spielman--Srivastava sparsification theorem (Kadison--Singer) gives \emph{edge}-light subsets with optimal constant.  The vertex-light variant addressed here is more delicate because correlations between edges sharing a vertex reduce the effective independence available for matrix Chernoff.
\item \textbf{Connection to spectral sparsification.}  An $\varepsilon$-light $S$ with $|S|\geq c\varepsilon n$ implies that $\bar S$ contains all but a $1-c\varepsilon$ fraction of vertices, with the boundary $L_{\partial S}+L_{V\setminus S}\succeq (1-\varepsilon)L$.  This is a one-sided spectral sparsification of $L$ supported on the vertex set $\bar S$.
\end{enumerate}

\end{document}
